%! Sinusoidal Function Transformations — Unit 6, Lesson 6.7 — Homework (Answer Key)
\documentclass[10pt]{article}
\usepackage{precalculus-article}
\usepackage{precalculus-key}   % pulls in -boxes; do NOT also load -boxes
\newcommand{\TallMath}[1]{$\displaystyle #1\rule[-1.4em]{0pt}{3.2em}$}

\begin{document}

\pageheader{Unit 6, Lesson 6.7}{Homework --- Answer Key}
\namedateperiod

\vspace{0.10in}

\begin{notesbox}{Problems 1--6 --- Key}

\textbf{1.}

\vspace{4pt}
\renewcommand{\arraystretch}{1.5}
\begin{tabular}{|l|c|c|}
\hline
\textbf{Function} & \textbf{Amplitude} & \textbf{Midline} \\
\hline
$f(\theta) = 7\sin(\theta) + 3$ & \ans{7} & \ans{$y=3$} \\
\hline
$g(\theta) = -5\cos(\theta) - 4$ & \ans{5} & \ans{$y=-4$} \\
\hline
$h(\theta) = \tfrac{3}{2}\sin(2\theta) + 1$ & \ans{$3/2$} & \ans{$y=1$} \\
\hline
$k(\theta) = -\cos(\pi\theta)$ & \ans{1} & \ans{$y=0$} \\
\hline
\end{tabular}

\vspace{0.12in}
\textbf{2.}

(a)\quad $b = $ \ans{4} \qquad Period $= $ \ans{$2\pi/4 = \pi/2$}

\vspace{6pt}
(b)\quad $b = $ \ans{$\pi/3$} \qquad Period $= \dfrac{2\pi}{\pi/3} = $ \ans{6}

\vspace{6pt}
(c)\quad Period $= \dfrac{2\pi}{1/4} = $ \ans{$8\pi$}

\vspace{0.12in}
\textbf{3.}

(a)\quad Factor: $2(\theta - \ans{$\pi/2$})$ \quad Phase shift: \ans{$\pi/2$ to the right}

\vspace{6pt}
(b)\quad Factor: $3(\theta + \ans{$\pi/6$})$ \quad Phase shift: \ans{$-\pi/6$ (left $\pi/6$ units)}

\vspace{6pt}
(c)\quad Factor: \ans{$\tfrac{1}{2}(\theta - \pi/2)$} \quad Phase shift: \ans{$\pi/2$ to the right}

\vspace{0.12in}
\textbf{4.}

(a)\quad Amplitude: \ans{4} \quad Period: \ans{$2\pi/3$} \quad
Phase shift: \ans{$\pi/6$ right} \quad Midline: \ans{$y=5$}

\vspace{8pt}
(b)\quad Factor: $4(\theta + \ans{$\pi/4$})$

Amplitude: \ans{2} \quad Period: \ans{$\pi/2$} \quad
Phase shift: \ans{$-\pi/4$ (left)} \quad Midline: \ans{$y=-3$}

\vspace{0.12in}
\textbf{5.} \ans{(B) $2\pi/3$}

Period $= 2\pi / |3| = 2\pi/3$. (The $-\pi$ and $+2$ do not affect the period.)

\vspace{0.12in}
\textbf{6.} \ans{(B) $f(\theta) = -6\sin(\theta) - 1$}

Amplitude $= |-6| = 6$. Midline $= -1$. (A) has midline $y=1$; (C) has midline $-1$ but max $=5\neq 5$... wait: (B) amplitude 6, midline $-1$. Correct.

\begin{teachernote}
Verify all options: (A) midline $y=1$, no. (B) amplitude $6$, midline $-1$, yes. (C) max $= -1+6=5$, min $= -1-6=-7$ --- midline $-1$ and amplitude 6, so (C) also works! However (C) has $d=-7$ and the formula shows midline $-7$, not $-1$. Re-check: $6\sin(\theta)-7$ has midline $y=-7$. So (B) is the only correct choice.
\end{teachernote}

\end{notesbox}

\vspace{0.08in}

\begin{extensionbox}
\textbf{Extension (Optional) --- Key}

(a)\quad Amplitude $= \dfrac{7-(-1)}{2} = $ \ans{4} \qquad
Midline $D = \dfrac{7+(-1)}{2} = $ \ans{3}

(b)\quad Period $= 4\pi$, so $\dfrac{2\pi}{B} = 4\pi \Rightarrow B = $ \ans{$1/2$}

(c)\quad $p(\theta) = 4\cos\!\left(\tfrac{1}{2}\theta\right)+3$.
Using $\cos\theta = \sin(\theta+\pi/2)$:
$p(\theta) = 4\sin\!\left(\tfrac{1}{2}\theta + \dfrac{\pi}{2}\right)+3
= 4\sin\!\left(\tfrac{1}{2}\!\left(\theta + \pi\right)\right)+3$

\ans{$p(\theta) = 4\sin\!\left(\tfrac{1}{2}(\theta+\pi)\right)+3$}
\end{extensionbox}

\vspace{0.08in}

\begin{tcolorbox}[colback=navylight, colframe=navy, arc=2mm,
    left=3mm, right=3mm, top=2mm, bottom=2mm]
\small\textbf{Preview --- Lesson 3.7: Sinusoidal Function Graphs}\\
You can now read amplitude, period, phase shift, and midline directly from
a sinusoidal equation. In Lesson 3.7 we'll go in both directions: from
equation to graph (using key points), and from a graph back to the equation.
Come ready to use today's parameters to plot and interpret complete sinusoidal graphs.
\end{tcolorbox}

\end{document}
